\documentclass[11pt]{article}

../../../macros/macros.tex
\lectureNoteRefs   % the lecture notes' section, exercise and lemma numbers, from refs/


% THIS FILE IS JUST FOR ENTERING YOUR NAME AND STUDENT ID.
% THERE IS NO PROBLEM 0.

% REPLACE THE TWO LINES BELOW AS INSTRUCTED

\setAuthor{
  \color{red} NAME (edit \texttt{problem\_0\_answer.tex)} % REPLACE THIS ENTIRE LINE WITH YOUR NAME
}

\setSID{%
  id% REPLACE THIS ENTIRE LINE WITH YOUR STUDENT ID
}


%%% Local Variables:
%%% mode: latex
%%% TeX-master: "homework_1"
%%% End:


\doHwHeader{Homework 1}

%%%%% EDIT THE FILES NAMED problem_X_answer.tex.  DON'T EDIT THIS FILE

%%%%%%%%%%%%%%%%%%%%%%%%%%%%%%%%%%%%%%%%%%%%%%%%%%%%%%%%%%%% 
%%%%%%%%%%%%%%%%%%%%%%%%%%%%%%%%%%%%%%%%%%%%%%%%%%%%%%%%%%%% 
%%%%%%%%%%%%%%%%%%%%%%%%%%%%%%%%%%%%%%%%%%%%%%%%%%%%%%%%%%%% 
%%%%%%%%%%%%%%%%%%%%%%%%%%%%%%%%%%%%%%%%%%%%%%%%%%%%%%%%%%%% 

\begin{document}


%%%%% EDIT THE FILES NAMED problem_X_answer.tex.  DON'T EDIT THIS FILE 

\paragraph{General instructions.} First read Lecture Notes 1, 2, and 3, including the exercises with answers,
then answer the problems given in the rest of this document.
Edit this LaTeX template to add your answers, as described in the file named \texttt{README.txt},
then compile it to produce \texttt{homework\_1.pdf}\ifGradescope{} for submission to gradescope\fi.

\paragraph{Don't know?  Say so.}
If you don't know the answer to a problem or any part of the problem,
we encourage you to simply write \emph{``I don't know''} instead of giving an answer.
(If you write ``I don't know'' you can also add any questions you might have.)

\paragraph{Gentle reminder: homework is for learning.}
As noted in Lecture Note 1, engaging fully with the homeworks is the best way to learn the material.
\emph{The primary role of homeworks is to help you learn the material,
  and the primary role of feedback on homeworks is to help you improve your ideas.}

Please try to explain your ideas as clearly as you can.
The clearer your explanations are, the more useful the feedback on them can be.

\paragraph{Collaboration and citation policy.}
If you are doing this homework as part of a course, follow the course's collaboration and citation policy.
In any case, at the end of each answer, list the people you collaborated with,
and cite every source that your answer uses ideas from
(other than the lecture notes and the sources listed in the lecture notes).
If there are none, write ``none''.

%%%%% EDIT THE FILES NAMED problem_X_answer.tex.  DON'T EDIT THIS FILE 

\vfill
\noindent{\footnotesize\copyrightNotice}

%%% Local Variables:
%%% mode: latex
%%% TeX-master: "homework_1"
%%% End:


%%%%%%%%%%%%%%%%%%%%%%%%%%%%%%%%%%%%%%%%%%%%%%%%%%%%%%%%%%%% 
%%%%%%%%%%%%%%%%%%%%%%%%%%%%%%%%%%%%%%%%%%%%%%%%%%%%%%%%%%%% 
%%%%%%%%%%%%%%%%%%%%%%%%%%%%%%%%%%%%%%%%%%%%%%%%%%%%%%%%%%%% 
%%%%%%%%%%%%%%%%%%%%%%%%%%%%%%%%%%%%%%%%%%%%%%%%%%%%%%%%%%%% 

%%%%% EDIT THE FILES NAMED problem_X_answer.tex.  DON'T EDIT THIS FILE 

%%%%%%%%%%%%%%%% PROBLEM 1  %%%%%%%%%%%%%%%%%

\newpage 

\problem \emph{(Section \ref{prelim: sec: background})}
Give the asymptotic values of the following functions, using the $\Theta$-notation.
Give the simplest possible answer (e.g., use $\Theta(n^5)$ instead of $\Theta(2n^5)$).
You don't have show your work, but you can if you want feedback on it
(and it may help you get partial credit if your answer is wrong). 
% 
% 
\begin{enumerate}[label=(\alph*)]
\item $\half n^3 + 7 n^4 - n^3 + 13n$
\item $2n^{-2} + 1/(\log^2 n)$
\item $n ( 7n^3\log^2 n + 9 n\log^5n) + 17^{\log_2 n}$
\item $n^2 2^n + 13n^4 + n^3\log n$
\item $n^5 4^n + n\, 5^n$
\end{enumerate}


% REPLACE THE SEVEN \REPLACEME PLACEHOLDERS BELOW WITH YOUR ANSWERS

\paragraph{Problem 1 answer}~

\begin{blue}

  \begin{enumerate}[label=(\alph*)]

  \item  % for $\half n^3 + 7 n^4 - n^3 + 13n$

    $\REPLACEME$
    
  \item  % for $2n^{-2} + 1/(\log^2 n)$

    $\REPLACEME$

  \item  % for $n ( 7n^3\log^2 n + 9 n\log^5n) + 17^{\log_2 n}$

    $\REPLACEME$
    
  \item % for $n^2 2^n + 13n^4 + n^3\log n$

    $\REPLACEME$
    
  \item % for $n^5 4^n + n\, 5^n$

    $\REPLACEME$

  \end{enumerate}

\end{blue}

\vfill

\paragraph{Problem \arabic{problem} collaboration and citations}~

I collaborated on this problem with the following people:
\begin{blue}
  \begin{itemize}
  \item \REPLACEME                    % write "none" if none
  \end{itemize}
\end{blue}

My answers use ideas from the following sources (other than the lecture notes and textbooks): 
\begin{blue}
  \begin{itemize}
  \item \REPLACEME                    % write "none" if none
  \end{itemize}
\end{blue}


%%% Local Variables:
%%% mode: latex
%%% TeX-master: "homework_1"
%%% End:


%%%%%%%%%%%%%%%%%%%%%%%%%%%%%%%%%%%%%%%%%%%%%%%%%%%%%%%%%%%% 
%%%%%%%%%%%%%%%%%%%%%%%%%%%%%%%%%%%%%%%%%%%%%%%%%%%%%%%%%%%% 
%%%%%%%%%%%%%%%%%%%%%%%%%%%%%%%%%%%%%%%%%%%%%%%%%%%%%%%%%%%% 
%%%%%%%%%%%%%%%%%%%%%%%%%%%%%%%%%%%%%%%%%%%%%%%%%%%%%%%%%%%% 

%%%%% EDIT THE FILES NAMED problem_X_answer.tex.  DON'T EDIT THIS FILE 

%%%%%%%%%%%%%%%% PROBLEM 2  %%%%%%%%%%%%%%%%%

\newpage 

\problem\emph{(Sections \ref{prelim: sec: background} and \ref{proofs: sec: exercises with solutions})}
Use the definition of big-O from Section \ref{prelim: sec: background} to prove the following theorem:

\begin{theorem}
  Suppose $f(n), g(n), h(n)$ are non-negative functions where $f(n) = O(h(n))$ and $g(n) = O(h(n))$.
  Then $f(n)\cdot g(n) = O(h^2(n))$.
\end{theorem}

Give the proof in long form (broken into small steps, just like the long-form proofs in the lecture notes). 
Be careful to get your reasoning about the quantifiers (in the definitions of big-O) right.


% REPLACE THE \REPLACEME PLACEHOLDERS BELOW WITH THE REMAINING STEPS IN YOUR PROOF

\paragraph{Problem 2 answer}~

\begin{longFormProof}

  \step Consider any functions $f$, $g$, and $h$ as described in the theorem.

  \begin{blue}
    
    \step \REPLACEME
    
    \step \REPLACEME

    % ETC... (YOU WILL NEED MORE THAN TWO STEPS)

  \end{blue}
  
  \step Therefore, by the definition of big-O, $f(n) \cdot g(n) = O(h^2(n))$. 
  
\end{longFormProof}

\vfill

% DON'T FORGET TO FILL OUT COLLABORATORS AND CITATIONS:

\paragraph{Problem \arabic{problem} collaboration and citations}~

I collaborated on this problem with the following people:
\begin{blue}
  \begin{itemize}
  \item \REPLACEME                    % write "none" if none
  \end{itemize}
\end{blue}

My answers use ideas from the following sources (other than the lecture notes and textbooks): 
\begin{blue}
  \begin{itemize}
  \item \REPLACEME                    % write "none" if none
  \end{itemize}
\end{blue}

%%% Local Variables:
%%% mode: latex
%%% TeX-master: "homework_1"
%%% End:



%%%%%%%%%%%%%%%%%%%%%%%%%%%%%%%%%%%%%%%%%%%%%%%%%%%%%%%%%%%% 
%%%%%%%%%%%%%%%%%%%%%%%%%%%%%%%%%%%%%%%%%%%%%%%%%%%%%%%%%%%% 
%%%%%%%%%%%%%%%%%%%%%%%%%%%%%%%%%%%%%%%%%%%%%%%%%%%%%%%%%%%% 
%%%%%%%%%%%%%%%%%%%%%%%%%%%%%%%%%%%%%%%%%%%%%%%%%%%%%%%%%%%% 

%%%%% EDIT THE FILES NAMED problem_X_answer.tex.  DON'T EDIT THIS FILE 

%%%%%%%%%%%%%%%% PROBLEM 3  %%%%%%%%%%%%%%%%%

\newpage

\problem\emph{(Induction)} {\small
  You encounter the following mysterious piece of code.

  \begin{myframed}
    \begin{steps}
    \item[] def $C(n)$:
      \step if  $n=0$: return $(0,1)$
      \step let $(u,w) = C(\lfloor n/2 \rfloor)$
      \step if $n$ is even:  return $(u+u, u + w)$
      \step if $n$ is odd: return $(u + w, w + w)$
    \end{steps}
  \end{myframed}

  \begin{enumerate}[label=(\alph*)] \small
  \item What are the outputs of $C(1)$, $C(2)$, $C(3)$, and $C(4)$?  Do not justify your answers.

  \item Give a long-form proof, by induction on $n$, that for all $n\ge 1$, executing $C(n)$ does at most
    $2 + 2 \cdot \log_2 n$ additions.  (Here we refer to the additions in Lines 3 and 4 of the algorithm.)

  \item Give a closed-form expression (in terms of $n$) for the value output by $C(n)$,
    assuming $n$ is a non-negative integer.
    Give a long-form proof, by induction on $n$, that $C(n)$ indeed outputs that value.

  \end{enumerate}
  % [review 2026-09-27] fix: Section 2.6.2's only inductive long-form proof is the deliberately wrong one (Exercise 2.9); point to a correct example too
  (For examples of inductive long-form proofs, see the proof in Section \ref{dc: sec: 1D local max} of \LN{dc},
  and Exercise \ref{proofs: exercise: understanding induction} in Section \ref{proofs: sec: exercises without solutions}, whose proof has the right form but contains an error for you to find.)
}


%%%%%%%%%% YOU KNOW THE DRILL... REPLACE ALL THE \REPLACEMEs WITH YOUR ANSWERS  %%%%%%%%%

\paragraph{Problem 3 answer}~

\begin{enumerate}[label=(\alph*)]

%%%%%%%%%%%%%%%%%%%%%%%%%%%%%%%%%%%%%%%%%%%%%%%%
%%%%%%%%%%%%%%%%%%%%%%%%%%%%%%%%%%%%%%%%%%%%%%%%

\item                           % (a)

  \(C(1) = \BLUE{
    \REPLACEME 
  }\) 

  \(C(2) = \BLUE{
    \REPLACEME 
  }\)

  \(C(3) = \BLUE{
    \REPLACEME 
  }\)

  \(C(4) = \BLUE{
    \REPLACEME 
  }\) 

%%%%%%%%%%%%%%%%%%%%%%%%%%%%%%%%%%%%%%%%%%%%%%%%
%%%%%%%%%%%%%%%%%%%%%%%%%%%%%%%%%%%%%%%%%%%%%%%%

\item                           % (b)

  \begin{lemma}
    For all $n\ge 1$, executing $C(n)$ does at most $2+2\log_2 n$ additions.
  \end{lemma}

  \begin{longFormProof}

    \step The proof is by induction on $n\ge 1$.

    \step\label{step: base additions}
    For the base case $n=1$, \BLUE{
      \REPLACEME 
    }

    \begin{block}\label{block: mysterious additions}

      \step For the inductive step, consider any $n\ge 2$, and assume the lemma holds for smaller values.

      \begin{blue}

        \step \REPLACEME  
 
        \step \REPLACEME

        % ETC. (YOU MAY NEED MORE THAN TWO STEPS.)

      \end{blue}
      
      \step So $C(n)$ does at most $2 + 2\log_2 n$ additions.
    \end{block}

    \step By Block~\ref{block: mysterious additions}, for any $n\ge 2$, if the lemma holds for smaller values,
    then it holds for $n$.

    \step By induction (with Step~\ref{step: base additions}), the lemma holds for all $n$.
    
  \end{longFormProof}


%%%%%%%%%%%%%%%%%%%%%%%%%%%%%%%%%%%%%%%%%%%%%%%%
%%%%%%%%%%%%%%%%%%%%%%%%%%%%%%%%%%%%%%%%%%%%%%%%

  \continued 

\item \emph{(Problem 3, continued)}   % part (c)
  
  \begin{lemma}
    For all $n\ge 0$, the value output by $C(n)$ is \(\BLUE{
      \REPLACEME
    }\)
  \end{lemma}
  
  \begin{longFormProof}
    
    \step The proof is by induction on $n$.

    \step\label{step: base}
    For the base case $n=0$, \BLUE{
      \REPLACEME
    }

    \begin{block}\label{block: mysterious}

      \step For the inductive step, consider any $n\ge 1$, and assume the lemma holds for smaller values.

      \step Consider the execution of $C(n)$.

      \begin{blue}
        
      \step \REPLACEME

      \step \REPLACEME

      % ETC. (YOU MAY NEED MORE THAN TWO STEPS)

      \end{blue}

      \begin{block}
        \step \underline{Case 1.} First consider the case that $n$ is odd.

        \begin{blue}
          
          \step \REPLACEME

          \step \REPLACEME

          % ETC. (YOU MAY NEED MORE THAN TWO STEPS)

        \end{blue}

      \end{block}

      \begin{block}
        \step \underline{Case 2.} Otherwise $n$ is even. 

        \begin{blue}

          \step \REPLACEME

          \step \REPLACEME

          % ETC. (YOU MAY NEED MORE THAN TWO STEPS) 

        \end{blue}
        
      \end{block}

      \step In either case, $C(n)$ returns \BLUE{\(
        \REPLACEME
        \).}
      That is, the lemma holds for this $n$.
    \end{block}

    \step By Block~\ref{block: mysterious}, for any $n\ge 1$, if the lemma holds for smaller values,
    then it holds for $n$.

    \step By induction (with Step~\ref{step: base}), the lemma holds for all $n$.
    
  \end{longFormProof}

\end{enumerate}


%%%%%%%%%%%%%%%%%%%%%%%%%%%%%%%%%%%%%%%%%%%%%%%%
%%%%%%%%%%%%%%%%%%%%%%%%%%%%%%%%%%%%%%%%%%%%%%%%

\vfill

\paragraph{Problem \arabic{problem} collaboration and citations}~

I collaborated on this problem with the following people:
\begin{blue}
  \begin{itemize}
  \item \REPLACEME                    % write "none" if none
  \end{itemize}
\end{blue}

My answers use ideas from the following sources (other than the lecture notes and textbooks): 
\begin{blue}
  \begin{itemize}
  \item \REPLACEME                    % write "none" if none
  \end{itemize}
\end{blue}


%%% Local Variables:
%%% mode: latex
%%% TeX-master: "homework_1"
%%% End:



%%%%%%%%%%%%%%%%%%%%%%%%%%%%%%%%%%%%%%%%%%%%%%%%%%%%%%%%%%%% 
%%%%%%%%%%%%%%%%%%%%%%%%%%%%%%%%%%%%%%%%%%%%%%%%%%%%%%%%%%%% 
%%%%%%%%%%%%%%%%%%%%%%%%%%%%%%%%%%%%%%%%%%%%%%%%%%%%%%%%%%%% 
%%%%%%%%%%%%%%%%%%%%%%%%%%%%%%%%%%%%%%%%%%%%%%%%%%%%%%%%%%%% 

%%%%% EDIT THE FILES NAMED problem_X_answer.tex.  DON'T EDIT THIS FILE 

%%%%%%%%%%%%%%%% PROBLEM 4  %%%%%%%%%%%%%%%%%

\newpage

\problem\emph{(Section \ref{matching: sec: problem definition})}

\begin{enumerate}[label=(\alph*)]

\item Define a specific instance $I$ of \prob{Stable Matching} with $n=3$
  that has only one stable matching.

\item Describe the stable matching.

\item Explain clearly why it is stable. 

\item Explain clearly why no other matching is stable.

\end{enumerate}


%%% YOU KNOW WHAT TO DO... %%%

\paragraph{Problem 4 answer}~

\begin{blue}

\begin{enumerate}[label=(\alph*)]

%%%%%%%%%%%%%%%%%%%%%%%%%%%%%%%%%%%%%%%%%%%%%%%%
%%%%%%%%%%%%%%%%%%%%%%%%%%%%%%%%%%%%%%%%%%%%%%%%

\item % (a) Define a specific instance $I$ of \prob{Stable Matching} with $n=3$ that has only one stable matching.
  
  \REPLACEME 

%%%%%%%%%%%%%%%%%%%%%%%%%%%%%%%%%%%%%%%%%%%%%%%%
%%%%%%%%%%%%%%%%%%%%%%%%%%%%%%%%%%%%%%%%%%%%%%%%

\item %  (b) Describe the stable matching.

  \REPLACEME 

%%%%%%%%%%%%%%%%%%%%%%%%%%%%%%%%%%%%%%%%%%%%%%%%
%%%%%%%%%%%%%%%%%%%%%%%%%%%%%%%%%%%%%%%%%%%%%%%%

\item % (c) Explain clearly why it is stable.

  \REPLACEME 

%%%%%%%%%%%%%%%%%%%%%%%%%%%%%%%%%%%%%%%%%%%%%%%%
%%%%%%%%%%%%%%%%%%%%%%%%%%%%%%%%%%%%%%%%%%%%%%%%

\item % (d) Explain clearly why no other matching is stable.
  
  \REPLACEME 

\end{enumerate}

\end{blue}

%%%%%%%%%%%%%%%%%%%%%%%%%%%%%%%%%%%%%%%%%%%%%%%%
%%%%%%%%%%%%%%%%%%%%%%%%%%%%%%%%%%%%%%%%%%%%%%%%

\vfill

\paragraph{Problem \arabic{problem} collaboration and citations}~

I collaborated on this problem with the following people:
\begin{blue}
  \begin{itemize}
  \item \REPLACEME                    % write "none" if none
  \end{itemize}
\end{blue}

My answers use ideas from the following sources (other than the lecture notes and textbooks): 
\begin{blue}
  \begin{itemize}
  \item \REPLACEME                    % write "none" if none
  \end{itemize}
\end{blue}


%%% Local Variables:
%%% mode: latex
%%% TeX-master: "homework_1"
%%% End:


%%%%%%%%%%%%%%%%%%%%%%%%%%%%%%%%%%%%%%%%%%%%%%%%%%%%%%%%%%%% 
%%%%%%%%%%%%%%%%%%%%%%%%%%%%%%%%%%%%%%%%%%%%%%%%%%%%%%%%%%%% 
%%%%%%%%%%%%%%%%%%%%%%%%%%%%%%%%%%%%%%%%%%%%%%%%%%%%%%%%%%%% 
%%%%%%%%%%%%%%%%%%%%%%%%%%%%%%%%%%%%%%%%%%%%%%%%%%%%%%%%%%%% 

%%%%% EDIT THE FILES NAMED problem_X_answer.tex.  DON'T EDIT THIS FILE 

%%%%%%%%%%%%%%%% PROBLEM 5 %%%%%%%%%%%%%%%%%

\newpage

\problem\emph{(Sections \ref{matching: sec: problem definition} and \ref{matching: sec: exercises with solutions}; \LN{proofs})} 
Recall the algorithm from Figure \ref{matching: fig: other alg}:

\begin{myframed}

  \begin{steps}
    \step For each pair $(i,j)$ in $\langle (1,1), (1,2), (2,1), (1,3), (3,1), (2,2), (2,3), (3,2), (3,3) \rangle$, do:

    \begin{steps}
      \step Consider those hospital/doctor pairs $(h, d)$ such that $d$ is $h$'s $i$th choice,
      $h$ is $d$'s $j$th choice, and neither $h$ nor $d$ are yet matched.
      For each such pair $(h, d)$, match $h$ to $d$.
    \end{steps}

    \step Return the matched pairs.
  \end{steps}

\end{myframed}

Consider the following proof.

\begin{lemma}
  For any execution of the algorithm, if it returns a perfect matching $M$, then $M$ is stable.
\end{lemma}
\begin{longFormProof}
  \step Consider any execution of the algorithm that returns a perfect matching $M$.

  \begin{block}\label{ZZZ}

    \step Consider an arbitrary hospital/doctor pair $(h, d)$.

    \begin{block}\label{WWW}
      \step Suppose for contradiction that $(h, d)$ is unstable for $M$.
      
      \step Let $d'$ and $h'$ be the respective partners of $h$ and $d$ in $M$ (they exist, as $M$ is perfect).

      \step Because $(h, d)$ is an unstable pair, $h$ ranks $d$ above $d'$.

      \step So, by inspection of the alg.~(Line 1), it considered the pair $(h, d)$ before the pair $(h, d')$. 

      \step\label{mnm} So, when the algorithm considered the pair $(h, d)$, hospital $h$ was not yet matched.

      \step Likewise, since $(h, d)$ is unstable for $M$, doctor $d$ ranks $h$ above $h'$.

      \step So, by inspection of the alg.~(Line 1), it considered the pair $(h, d)$ before the pair $(h', d)$. 

      \step\label{wnm} So, when the algorithm considered the pair $(h, d)$, doctor $d$ was not yet matched.

      \step By Steps~\ref{mnm} and~\ref{wnm},
      when the alg.~considered $(h, d)$, neither $h$ nor $d$ were matched.

      \step So, by inspection of the algorithm, it must have matched $h$ to $d$ in that iteration.

      \step But once a pair is matched, it stays matched, so $(h, d)$ is in $M$.

      \step This is a contradiction, as no matched pair is unstable.
    \end{block}

    \step By Block~\ref{WWW}, $(h, d)$ is not unstable for $M$. 
  \end{block}

  \step By Block~\ref{ZZZ}, no pair is unstable for $M$.  That is, $M$ is stable.
\end{longFormProof}

\begin{enumerate}[label=(\alph*)]

\item
  What is the first line in this proof that makes an assertion that is not necessarily true?

\item
  What is the logical error at that point in the reasoning?

\item
  Describe an input instance on which the algorithm fails.
  Use hospitals $X, Y, Z$ and doctors $A, B, C$.
  Describe the instance by specifying the $3\times 3$ table of mutual ranks for each pair,
  as in the solution to Exercise \ref{matching: prob: before gale shapley}.

\item
  Describe the execution of the algorithm on that input, and the resulting matching.

\item
  What is the resulting matching?

\item
  What is a pair $(h, d)$ that is unstable for the matching?  Why is it unstable?
  
\end{enumerate}

\continued

Here is the algorithm for reference:

\begin{myframed}

  \begin{steps}
    \step For each pair $(i,j)$ in $\langle (1,1), (1,2), (2,1), (1,3), (3,1), (2,2), (2,3), (3,2), (3,3) \rangle$, do:

    \begin{steps}
      \step Consider those hospital/doctor pairs $(h, d)$ such that $d$ is $h$'s $i$th choice,
      $h$ is $d$'s $j$th choice, and neither $h$ nor $d$ are yet matched.
      For each such pair $(h, d)$, match $h$ to $d$.
    \end{steps}

    \step Return the matched pairs.
  \end{steps}

\end{myframed}

%%%%%%%%%%%%%%%%%%%%%%%%%%%%%%%%%%%%%%%%%%%%%%%%%%%%%%%%%%%% 
%%%%%%%%%%%%%%%%%%%%%%%%%%%%%%%%%%%%%%%%%%%%%%%%%%%%%%%%%%%% 

%%%%% EDIT THE FILES NAMED problem_X_answer.tex.  DON'T EDIT THIS FILE 

\paragraph{Problem 5 answer}~

\begin{enumerate}[label=(\alph*)]

%%%%%%%%%%%%%%%%%%%%%%%%%%%%%%%%%%%%%%%%%%%%%%%%%%%%%%
%%%%%%%%%%%%%%%%%%%%%%%%%%%%%%%%%%%%%%%%%%%%%%%%%%%%%%

\item % (a) What is the first line in this proof that makes an assertion that is not necessarily true?

  \BLUE{
    \REPLACEME
  }

%%%%%%%%%%%%%%%%%%%%%%%%%%%%%%%%%%%%%%%%%%%%%%%%%%%%%%
%%%%%%%%%%%%%%%%%%%%%%%%%%%%%%%%%%%%%%%%%%%%%%%%%%%%%%

\item % (b) What is the logical error at that point in the reasoning?

  \BLUE{
    \REPLACEME
  }

%%%%%%%%%%%%%%%%%%%%%%%%%%%%%%%%%%%%%%%%%%%%%%%%%%%%%%
%%%%%%%%%%%%%%%%%%%%%%%%%%%%%%%%%%%%%%%%%%%%%%%%%%%%%%

\item % (c)
  % Describe an input instance on which the algorithm fails.
  % [review 2026-09-27] fix: comment said men/women; the problem uses hospitals X, Y, Z and doctors A, B, C
  % Use hospitals $X, Y, Z$ and doctors $A, B, C$.
  % Describe the instance by specifying the $3\times 3$ table of mutual ranks for each pair,
  % as in the solution to Exercise 3.4.

  Here is the input, specified by each pair's mutual ranking:
  
  % REPLACE THE TABLE ENTRIES BELOW BY THE MUTUAL RANKS FOR YOUR INSTANCE:
  
  \begin{blue}
    \[
    \begin{array}{c|ccc}
       & A & B & C \\\hline
      X & (0, 0) & (0, 0) & (0, 0)  \\
      Y & (0, 0) & (0, 0) & (0, 0)  \\
      Z & (0, 0) & (0, 0) & (0, 0)
    \end{array}
    \]
  \end{blue}

%%%%%%%%%%%%%%%%%%%%%%%%%%%%%%%%%%%%%%%%%%%%%%%%%%%%%%
%%%%%%%%%%%%%%%%%%%%%%%%%%%%%%%%%%%%%%%%%%%%%%%%%%%%%%

\item %  (d) Describe the execution of the algorithm on that input, and the resulting matching.

  Here is what happens during the execution of the algorithm on that input:

  \begin{blue}

    \REPLACEME

  \end{blue}

%%%%%%%%%%%%%%%%%%%%%%%%%%%%%%%%%%%%%%%%%%%%%%%%%%%%%%
%%%%%%%%%%%%%%%%%%%%%%%%%%%%%%%%%%%%%%%%%%%%%%%%%%%%%%

\item % (e) What is the resulting matching?

  Here is the resulting matching:

  \begin{blue}
    \( \REPLACEME \)
  \end{blue}
  
%%%%%%%%%%%%%%%%%%%%%%%%%%%%%%%%%%%%%%%%%%%%%%%%%%%%%%
%%%%%%%%%%%%%%%%%%%%%%%%%%%%%%%%%%%%%%%%%%%%%%%%%%%%%%

\item %  (f) What is a pair $(m, w)$ that is unstable for the matching?  Why is it unstable? 
  
  The pair \BLUE{
    \(\REPLACEME\)
  } is unstable, because \BLUE{
    \REPLACEME
  }
  
\end{enumerate}

%%%%%%%%%%%%%%%%%%%%%%%%%%%%%%%%%%%%%%%%%%%%%%%%%%%%%%
%%%%%%%%%%%%%%%%%%%%%%%%%%%%%%%%%%%%%%%%%%%%%%%%%%%%%%

\vfill

\paragraph{Problem \arabic{problem} collaboration and citations}~

I collaborated on this problem with the following people:
\begin{blue}
  \begin{itemize}
  \item \REPLACEME                    % write "none" if none
  \end{itemize}
\end{blue}


My answers use ideas from the following sources (other than the lecture notes and textbooks): 
\begin{blue}
  \begin{itemize}
  \item \REPLACEME                    % write "none" if none
  \end{itemize}
\end{blue}
 
%%% Local Variables:
%%% mode: latex
%%% TeX-master: "homework_1"
%%% End:



\end{document}

%%% Local Variables:
%%% mode: latex
%%% TeX-master: t
%%% End:
